\chapter{结论与展望}\label{chap:conclusion}

\section{结论}

针对异常检测评测中的跨模态割裂与标签污染鲁棒性缺失两大问题，本文构建了一个覆盖表格、时序、图、图像与文本五大模态的统一抗污染评测基准，在 29 个真实数据集上系统评测了 26 种算法，并设计了一套即插即用的主动去噪防御。主要结论如下：

\begin{enumerate}
    \item \textbf{统一评测底座}。我们实现了可对齐五大模态、以一致接口驱动 26 种算法的自动化评测流水线，使不同算法族群得以在相同的污染标准下横向比较。
    \item \textbf{噪声退化规律}。在干净标注下监督模型占优（Exp1）；随污染加剧，其退化幅度高度依赖模态，在文本、时序上跌幅可达 16--20 个百分点，而无监督方法因标签免疫而保持稳定，并在时序模态出现监督/无监督性能反转（Exp2、Exp3）。
    \item \textbf{主动防御的有效性与边界}。基于样本剪裁的 \texttt{IQR\_Trim} 在 20\% 污染下使 MLP 的 AUC-ROC 回升 2.52 个百分点，且剪裁策略稳定优于标签纠正；但其对逻辑回归与 TabPFN 产生负迁移，揭示了去噪防御的失效边界与“没有免费的午餐”原则（Exp4）。
\end{enumerate}

\section{展望}

本文仍有若干可拓展的方向：
\begin{enumerate}
    \item \textbf{更真实的噪声模型}。本文聚焦对称标签翻转，未来可研究非对称与特征依赖的标签噪声（即靠近边界的模糊样本更易被误标），并结合置信学习设计自适应去噪。
    \item \textbf{大模型辅助的可解释异常检测}。可探索以推理大模型 \cite{guo_deepseek-r1_2025} 对时序、文本日志与图属性异常进行零样本判定与自然语言解释。
    \item \textbf{联邦式主动防御}。在隐私保护的多端场景下，研究不汇聚原始数据的分布式去噪与协同净化机制。
\end{enumerate}
